% =====================================================================
%  Выводы по лабораторной работе.
% =====================================================================

\labpart{Выводы}

В ходе выполнения лабораторной работы реализована система
автоматического распознавания языка текста по варианту 25 (английский и
французский языки, HTML-документы, методы частотных слов, алфавитный и
нейросетевой) как новый сценарий использования поверх архитектуры,
разработанной в ЛР1: добавлен один новый сервис (\texttt{lang-id-service}),
пять новых таблиц БД и новая страница интерфейса, при этом ни один
существующий контракт системы (поиск, краулинг, метрики поиска) не
был нарушен. Все три метода подчиняются единой трёхшаговой схеме
(построение ПОЯ по обучающему корпусу, построение ПОД для входного
документа, выбор языка с минимальным расстоянием между ними), различаясь
только правилом построения профиля и метрикой расстояния, что позволило
реализовать общий, покрытый тестами интерфейс сравнения методов
(\texttt{app.domain.lang\_id}).

На собранной тестовой коллекции \texttt{Wikipedia} (1000 статей,
поровну с \texttt{en.wikipedia.org} и \texttt{fr.wikipedia.org},
800 обучающих / 200 тестовых документов) алфавитный и нейросетевой
методы показали 100\,\% точность, а метод частотных слов -- 95{,}5\,\%
(табл.~\ref{tab:metrics-langid}) -- ожидаемый результат для языковой
пары со столь разными системами письма и словарным составом. По
быстродействию нейросетевой метод оказался почти в три раза быстрее
обоих лексических (1{,}24~мс против 3{,}96--4{,}97~мс на документ) за
счёт того, что классификация выполняется одним прямым проходом над уже
готовым эмбеддингом документа, без повторного анализа всего текста.
Единственная методологическая оговорка -- точность 100\,\% у двух из
трёх методов может отчасти объясняться тем, что использовалось одно
фиксированное стратифицированное разбиение на обучающую и тестовую
выборки, а не усреднение по нескольким случайным разбиениям
(кросс-валидация); на менее контрастной паре языков или при другом
разбиении разрыв между методами, вероятно, был бы заметнее. Разработанный
интерфейс поддерживает как пакетное тестирование на всей размеченной
коллекции, так и разовую проверку произвольной веб-страницы или текста,
с выгрузкой результата в CSV/JSON и печатью, что закрывает
соответствующее требование методических указаний.
